\section{Related work}
\label{sec:related}

This paper sits in a literature that has established most of its framing, and the honest way to
position it is to say what each neighbouring work did before saying what is left. We group the
work by theme and end the section with the boundary stated exactly.

\subsection{AI and the public finances}

The proposition that labor-displacing AI is a fiscal event before it is anything else is
established, and this paper does not claim it.

\citet{PriceSuresh2026} measure it directly. Working in a budget model, they put about
eighty-four percent of federal revenue on individual or payroll taxes and roughly two thirds of it
directly on labor, and run four AI substitution scenarios on two axes, one of which is the
reemployment rate and the other the pricing of AI services. Those are the same two quantities this
paper calls the retained wage share and the effective rate on AI capital. Their policy options are
explicitly illustrative and unsized. Because output is held constant by construction, their
exercise is the analogue of our output-preserved case alone, and it carries no household debt, no
mortgages and no bank balance sheets.

\citet{Barhoumi2026} is a workshop synthesis rather than a model, and it contains no numbers; what
matters for our purposes is that it states the mechanism this paper measures. It says directly
that large-scale job losses weaken household balance sheets and raise default risks in banking
systems exposed to consumer credit, and it states the distributional asymmetry, that the costs
fall on incumbent workers and borrowers while the gains accrue to new entrants and AI-intensive
firms. Both sides of the exposure appear in the same paragraph of that note. What is absent is any
measurement of them, and any map of who holds them. It is the closest approach to this paper's
question that we located, and we say so rather than leaving a reader to find it.

\citet{IeongSaputraManiarCheng2026} build a four-channel account of how a labor-displacing
transition erodes a tax base, simulated for an average OECD country: the shift from
highly-taxed labor income to lower-taxed capital income, leakage when AI providers are
non-resident, consumer benefits that are never taxed, and higher social support spending. The
first channel is the fiscal condition of Section~\ref{sec:fiscal} stated in a different accounting
frame. The scenarios in the workshop synthesis above were developed in collaboration with this
group, which connects two works that might otherwise look independent, and we treat all three as
connected work rather than as competition.

\citet{KorinekLockwood2026} ask the optimal-taxation question across stages of an AI transition.
That is a different question from ours. We measure what the existing tax system does reach; they
characterise what an optimal one should do. \citet{FalkTsoukalas2026} make an argument that is
stronger than ours and reaches it first: each firm captures the full cost saving from automation
while bearing only a fraction of the demand loss it creates, so capital income taxes cannot
resolve the externality and a Pigouvian automation tax is the instrument that can. Our
contribution here is not the verdict but the calibration, and our measured effective rate can be
read as an input to their argument rather than a competitor to it.
\citet{CasasTorres2024} established the base-shift mechanism two years before this project began,
showing in general equilibrium that as automation rises the ratio of fiscal revenue to output
falls, because the inputs that bear the tax are replaced by an input that does not.
\citet{Brollo2024} supplies the one measured policy effect we rely on in
Section~\ref{sec:institutions}, that more generous unemployment insurance substantially attenuates
the wage decline from robotisation.

\paragraph{What this paper does differently.} None of these decomposes the stock of financial
claims by the income that services it, and none records who holds each side. They establish that
the revenue falls. We measure whose balance sheets the claims sit on when it does.

\subsection{Automation, labor demand and reemployment}

\citet{AcemogluRestrepo2020} is the empirical anchor for what displacement does to employment and
wages, and the source of the reemployment structure our retained wage share is built on.
\citet{AcemogluManeraRestrepo2020} does double duty here. It supplies the tax result that
Section~\ref{sec:fiscal} takes as derived rather than assumed, that under full expensing the
normal return is relieved at the entity level so the marginal rate on it is the owner's rate
alone, and it is the source of the observation that the US tax system taxes capital more lightly
than labor at the margin. \citet{Huckfeldt2022} gives the wage loss on reemployment that our
retained wage share uses, and the finding that the loss is concentrated among occupation switchers
rather than spread evenly.

\citet{ManningAguirreMuroMethkupally2026} build an occupation-level index of adaptive capacity and
find, against what one might expect, that AI exposure and adaptive capacity are positively
correlated: most workers in the most exposed occupations are in occupations with above-median
capacity to absorb a transition. That is the same conclusion this paper reaches from household
microdata in Section~\ref{sec:displacement}, on a different object: exposure alone does not
identify vulnerability, and pay does most of the work.

\paragraph{What this paper does differently.} An occupation-level index cannot show that the
relationship between buffers and exposure runs through pay level and disappears once pay is
controlled, because that is a household-level result requiring household balances. Nor can it
show that the thinnest buffers sit in the second wage quintile rather than the first.

\subsection{Capital taxation and the ownership of corporate claims}

Our effective rate is assembled from a literature that has measured its components separately.
\citet{Auerbach2017} sets out the structure of taxing the corporate return and the sense in which
expensing exempts the normal return. \citet{Barkai2020} supplies one of the two rent share
readings we carry, and the incompatibility between that reading and its main alternative is why we
report both and never average them. \citet{TorslovWierZucman2023} supplies the profit-shifting
share, which we report because it matters and rank fourth among seven levers because it does not
matter most. The ownership parameters that turn out to be decisive, the share of corporate equity
in taxable accounts, the deferral of tax between accrual and realization, the share of gains never
taxed because basis is stepped up at death, and the tax status of corporate debt, come from work
in the public finance and budget-office literature that we cite at the point of use.

\paragraph{What this paper does differently.} These components have not, so far as we located,
been assembled into a single marginal effective rate on the income that moves from wages to AI
capital, and compared against the rate that revenue replacement would require. The assembly is
where the ownership explanation becomes visible: the reason the rate is low is who owns the
claims, not where the profits are booked.

\subsection{Household balance sheets, default and stress testing}

The household side of this paper rests on a literature about what job loss does to a borrower.
\citet{BhuttaBrickerDettlingKelliherLaufer2019} and \citet{MerikullRoom2020} establish the
relationship between income shocks, liquid buffers and distress that our engine's uplifts rely on,
and \citet{AlbaceteFessler2010} is the stress-testing template our engine follows in structure. Our
own default uplifts are taken from the literature on job loss and mortgage default rather than
estimated here. \citet{CohenKillenLau2026} supplies the measured size of bank commitments to
AI-adjacent industries that Section~\ref{sec:bust} uses, and is the one source in this paper that
sits on both sides of the exposure at once.

\citet{Chen2026} is the largest overlap with this paper, and it is with our second round. It is a
macro-financial stress test of rapid AI adoption with eleven testable predictions, reaching private
credit and mortgages, and it makes the point that AI exposure is concentrated among high-income
households who account for a large share of consumption. What remains ours is severity measured
against a named supervisory benchmark, loss rates by loan category, application across individual
balance sheets, and the decomposition of the loss by holder.

\paragraph{The closest existing statistic.} The Federal Reserve's debt service ratio measures
household debt service against all household income, and the distributional financial accounts
give balance sheets by wealth percentile. Both are closer to our object than anything else
published. Neither decomposes the servicing cash flow by the income that supplies it, and neither
extends beyond households to the rest of the claim structure.

\subsection{What is ours}

\begin{quote}
The fiscal mechanism is established. What this paper adds is the measured holder structure of the
claims that wages pay, on both sides of the exposure and on the same basis; household losses from
survey microdata mapped to the supervisor's own loss rates; the incidence of those losses by who
bears them; a capital tax condition assembled from sourced components and independently rebuilt;
and a public claims register in which every reported quantity carries its status and its
replication flag.
\end{quote}

One boundary statement belongs here rather than in the limitations, because it qualifies the
contribution rather than the result. Our search located no prior construction of the labor backing
accounts as an object. That search was not systematic, so the claim we make is ``none located''
rather than ``none exists''. We have already retired one priority claim in this project after a
non-systematic search found prior work, and the same could happen here.
